\section{Experimental Results}
\label{sec:results}

Every numerical result in this section is linked deterministically via LaTeX macros from our \TotalSealedRuns{} sealed experiment suites (\TotalSealedRecords{} total runs, \TotalEscapedDefects{} escaped defects).

% AUTO-GENERATED FROM v4_exp2_main_bench — DO NOT EDIT
\begin{table*}[t]
\centering
\footnotesize
\setlength{\tabcolsep}{4.2pt}
\caption{Main Benchmark Across 3 Real OSS Traces (450 Eval PRs, $5\times$ Velocity, $\phi=5\%$, $\rho=80\%$, 20 Paired Seeds $\times$ 3 Repos; 0 Cumulative Escapes. Throughput, P95, and Cost Show Mean [95\% Marginal Bootstrap CI]; Rej.\ Shows Mean Defective/Persistent/Flaky Counts).}
\label{tab:main_benchmark}
\begin{tabular}{lcccccc}
\toprule
\textbf{Strategy} & \textbf{Throughput (PR/h)} $\uparrow$ & \textbf{P95 Lat.\ (min)} $\downarrow$ & \textbf{CI Cost (m/PR)} $\downarrow$ & \textbf{Inv.\ (\%)} $\downarrow$ & \textbf{FR (\%)} $\downarrow$ & \textbf{Rej.\ (D/P/F)} \\
\midrule
SerialQueue & 3.75 [3.68, 3.81] & 1413.9 [1406.2, 1421.3] & 16.11 [15.80, 16.43] & 0.0\% & 4.36\% & 11.8/19.8/6.0 \\
StaticBatchBisect & 16.92 [16.32, 17.51] & 113.8 [105.6, 122.4] & 18.25 [17.35, 19.22] & 0.9\% & 2.01\% & 11.8/20.2/2.8 \\
SpeculativeLinear & 5.82 [5.57, 6.06] & 846.8 [818.6, 875.0] & 28.63 [27.44, 29.86] & 38.3\% & 2.12\% & 11.8/20.0/2.9 \\
FIFOOutcomeTree & 11.05 [10.59, 11.50] & 319.4 [304.7, 334.4] & 39.90 [38.35, 41.45] & 75.2\% & 1.78\% & 11.8/20.0/2.5 \\
DepAwareOutcomeTree & 11.75 [11.19, 12.29] & 283.5 [265.1, 302.0] & 36.95 [35.17, 38.87] & 71.5\% & 2.19\% & 11.8/20.0/3.0 \\
FlakeAwareBatch & 13.44 [12.79, 14.09] & 215.8 [197.5, 234.4] & 27.34 [25.66, 29.06] & 1.9\% & 0.12\% & 11.8/20.6/0.2 \\
\textbf{SSMQ} & \textbf{18.18 [17.66, 18.69]} & \textbf{81.0 [75.4, 86.7]} & \textbf{16.75 [16.02, 17.51]} & \textbf{2.0\%} & \textbf{1.17\%} & \textbf{11.8/20.2/1.6} \\
\bottomrule
\end{tabular}
\end{table*}


\subsection{RQ1: End-to-End Throughput, Latency, and CI Compute}
Table~\ref{tab:main_benchmark} summarizes aggregate performance across all three open-source repositories at $5\times$ AI-agent PR velocity ($30$ PRs/hr) and $5\%$ test flakiness. \textbf{SSMQ} achieves the strongest performance across all primary dimensions simultaneously:
\begin{itemize}
    \item \textbf{Throughput}: SSMQ merges \SSMQThroughput{} PRs/hr, outperforming \texttt{StaticBatchBisect} (\StaticThroughput{} PRs/hr, +\GainThroughputVsStaticPct{}\%), \texttt{FlakeAwareBatch} (\FlakeBatchThroughput{} PRs/hr, +\GainThroughputVsFlakePct{}\%), \texttt{SpeculativeLinear} (\SpecLinThroughput{} PRs/hr, +\GainThroughputVsSpecLinPct{}\%), \texttt{UberSubmitQueue} (\UberThroughput{} PRs/hr, +\GainThroughputVsUberPct{}\%), and \texttt{SerialQueue} (\SerialThroughput{} PRs/hr).
    \item \textbf{Tail Merge Latency}: SSMQ compresses P95 merge latency to \SSMQPNinetyFiveLat{}~min (P50 \SSMQPFiftyLat{}~min), representing a \LatRedVsStaticPct{}\% reduction over \texttt{StaticBatchBisect} (\StaticPNinetyFiveLat{}~min), \LatRedVsSpecLinPct{}\% over \texttt{SpeculativeLinear} (\SpecLinPNinetyFiveLat{}~min), and \LatRedVsUberPct{}\% over \texttt{UberSubmitQueue} (\UberPNinetyFiveLat{}~min; Cliff's $|\delta|=\CliffsDeltaLatVsUber{}$).
    \item \textbf{CI Compute Efficiency}: SSMQ consumes only \SSMQCICost{} CI runner-minutes per merged PR---saving \CostSaveVsStaticPct{}\% vs.\ \texttt{StaticBatchBisect} (\StaticCICost{}~min/PR, Cliff's $|\delta|=\CliffsDeltaCostVsStatic{}$), \CostSaveVsUberPct{}\% vs.\ \texttt{UberSubmitQueue} (\UberCICost{}~min/PR, $|\delta|=\CliffsDeltaCostVsUber{}$), \CostSaveVsFlakePct{}\% vs.\ \texttt{FlakeAwareBatch} (\FlakeBatchCICost{}~min/PR), and \CostSaveVsSpecLinPct{}\% vs.\ \texttt{SpeculativeLinear} (\SpecLinCICost{}~min/PR).
    \item \textbf{Invalidation & Fault Isolation}: Linear speculative trains suffer a \SpecLinInvalidRatePct{}\% speculative invalidation rate due to cascading downstream aborts, whereas SSMQ limits invalidation to \SSMQInvalidRatePct{}\% and reduces mean blame rounds per faulty batch from \StaticBlameRounds{} (\texttt{StaticBatchBisect}) and \FlakeBatchBlameRounds{} (\texttt{FlakeAwareBatch}) down to \SSMQBlameRounds{} (-\BlameRedVsStaticPct{}\% vs.\ static bisection).
\end{itemize}

\textbf{Velocity Surge Scaling ($1\times$ to $10\times$).} Under extreme $10\times$ arrival velocity (\texttt{exp2\_velocity\_sweep\_v1}), linear and risk-gated single-train queues hit a hard saturation ceiling (\texttt{UberSubmitQueue} stalls at \UberTenXThroughput{} PRs/hr with \UberTenXPNinetyFiveLat{}~min P95 latency). In contrast, SSMQ scales smoothly to \SSMQTenXThroughput{} PRs/hr (vs.\ \StaticTenXThroughput{} PRs/hr for \texttt{StaticBatchBisect}) while keeping P95 latency bounded at \SSMQTenXPNinetyFiveLat{}~min.

% AUTO-GENERATED FROM v4_exp4_ablation — DO NOT EDIT
\begin{table}[t]
\centering
\footnotesize
\setlength{\tabcolsep}{3.0pt}
\caption{Deconfounded Component Ablation ($5\times$, $\phi=5\%$, $\rho=80\%$, 20 Seeds $\times$ 3 Repos; RootSz = Mean Root Batch Size).}
\label{tab:ablation}
\begin{tabular}{lccccc}
\toprule
\textbf{Variant} & \textbf{Tput} & \textbf{P95} & \textbf{Cost} & \textbf{FR\%} & \textbf{RootSz} \\
\midrule
\textbf{SSMQ} & \textbf{18.18} & \textbf{81.0} & \textbf{16.75} & \textbf{1.17\%} & \textbf{1.44} \\
SSMQ-NoCone & 17.15 & 102.2 & 17.98 & 0.92\% & 2.09 \\
SSMQ-NoRisk & 18.09 & 82.3 & 16.80 & 1.11\% & 1.45 \\
SSMQ-NoGraph & 16.41 & 125.7 & 19.47 & 0.79\% & 2.80 \\
SSMQ-NoFlaky & 18.99 & 65.4 & 15.06 & 2.63\% & 1.43 \\
SSMQ-FixedRetry & 16.38 & 113.8 & 19.76 & 0.13\% & 1.45 \\
SSMQ-NoBlame & 18.09 & 81.8 & 16.81 & 1.21\% & 1.44 \\
StaticBatch & 16.92 & 113.8 & 18.25 & 2.01\% & 2.72 \\
\bottomrule
\end{tabular}
\end{table}


\subsection{RQ2: Component Ablation Analysis}
Table~\ref{tab:ablation} isolates each SSMQ mechanism under identical paired traces (\texttt{exp3\_ablation\_v1}):
\begin{enumerate}
    \item \textbf{Graph-Disjoint Batching is the Dominant Lever}: Removing dependency-graph packing (\texttt{SSMQ-NoGraph}) degrades P95 latency from \SSMQPNinetyFiveLat{}~min to \AblNoGraphLat{}~min and increases CI compute from \SSMQCICost{} to \AblNoGraphCICost{}~min/PR. Without graph-disjoint lanes, FIFO batching bundles overlapping PRs that trigger avoidable semantic conflicts.
    \item \textbf{1-Step Causal Blame Reduces Recovery Overhead}: Disabling causal blame (\texttt{SSMQ-NoBlame}) increases mean blame rounds from \SSMQBlameRounds{} to \AblNoBlameRounds{} and nudges CI compute from \SSMQCICost{} to \AblNoBlameCICost{}~min/PR. Because graph-disjoint batching already eliminates intra-batch semantic conflicts before CI runs, blame attribution activates mainly on true single-PR defects; its marginal impact grows substantially whenever conflict rates or batch sizes increase.
    \item \textbf{Flaky-Rerun Trade-Off Boundary}: Interestingly, disabling selective flaky reruns (\texttt{SSMQ-NoFlaky}) yields slightly lower P95 latency (\AblNoFlakyLat{}~min) and CI cost (\AblNoFlakyCICost{}~min/PR) at $\phi=5\%$. This reveals an important empirical trade-off: when graph-disjoint batches already have high pass rates and moderate batch size, rerunning a failed batch once incurs a full extra build on genuinely defective batches, whereas splitting/blaming immediately isolates true defects faster. Conversely, unconditional reruns (\texttt{FlakeAwareBatch}) severely inflate both latency (\FlakeBatchPNinetyFiveLat{}~min) and cost (\FlakeBatchCICost{}~min/PR).
\end{enumerate}

\subsection{RQ3: Robustness Under Extreme Flakiness and Imperfect Blame}
\textbf{Flaky Test Sweep ($0\%$--$10\%$).} Across \texttt{exp4a\_flaky\_sweep\_v1}, as non-deterministic flakiness rises to an extreme $\phi=10\%$, \texttt{SpeculativeLinear} degrades sharply to \SpecLinFlakyTenCost{} CI min/PR and \texttt{UberSubmitQueue} reaches \UberFlakyTenCost{}~min/PR, whereas SSMQ remains the most efficient scheduler at \SSMQFlakyTenCost{}~min/PR.

\textbf{Imperfect Causal Attribution ($50\%$--$95\%$).} Across \texttt{exp4b\_blame\_accuracy\_v1}, varying causal blame Top-1 accuracy from $95\%$ down to $50\%$ increases mean blame rounds only slightly from \BlameNinetyFiveRounds{} to \BlameFiftyRounds{}---still less than half the \StaticBlameRounds{} rounds required by \texttt{StaticBatchBisect}. Because SSMQ validates evicted suspects in parallel and falls back deterministically to binary bisection on misattribution, imperfect LLM-based fault localization never compromises mainline safety or collapses queue performance.
